\section{Conclusion}

In this work, we studied hallucinated detections in object detectors under out-of-distribution (OoD) settings, focusing on the role of frequency bias and object scale in objectness miscalibration. We proposed a lightweight, frequency-aware calibration method that works at inference time without requiring additional training or model changes.

Our experiments show that the proposed approach reduces false positives and hallucination rates while maintaining in-distribution performance. In addition, explainability analyses using Grad-CAM, occlusion sensitivity, and insertion-deletion metrics indicate improved alignment between model attention and meaningful object regions.

Overall, our results highlight the importance of frequency-based cues in understanding and improving OoD robustness. By combining OoD mitigation with explainability, this work provides both practical improvements and insights into model behavior under distribution shift.

Future work might include extending the method to other architectures and datasets, as well as exploring more causal explainability techniques and using robust and faster models to further enhance performance and interpretability.